% Chương 3 — nội dung IR của M2. Dùng bằng % Chương 3 — nội dung IR của M2. Dùng bằng % Chương 3 — nội dung IR của M2. Dùng bằng % Chương 3 — nội dung IR của M2. Dùng bằng \input{chapter3_ir} trong main.tex.
% Chỉ dùng các gói đã khai báo trong main.tex: report, vietnamese, amsmath,
% amssymb, graphicx, hyperref, geometry. Không thêm gói mới.
\chapter{THIẾT KẾ HỆ THỐNG TÌM KIẾM VĂN BẢN}

\section{Phạm vi trách nhiệm của M2 trong kiến trúc}

Kiến trúc tổng thể chia thành ba lớp trách nhiệm. M1 sở hữu tiền xử lý ngôn ngữ
tự nhiên, M2 sở hữu biến các từ đã tách thành dữ liệu lập chỉ mục và truy xuất,
M3 sở hữu giao diện người dùng. Chương này trình bày phần của M2.

Các lớp và ranh giới tích hợp được trình bày trong Bảng~\ref{tab:kien-truc-m2}.
\begin{table}[ht]
\centering
\begin{tabular}{|p{0.08\textwidth}|p{0.43\textwidth}|p{0.35\textwidth}|}
\hline
Lớp & Thành phần & File chính \\
\hline
M1 & Maximum Matching, stopword tầng 1 & \texttt{engine/tokenizer/} \\
M2 & Adapter M1 có fingerprint & \texttt{engine/preprocessing.py} \\
M2 & TDV, TF-IDF, chỉ mục đảo, tìm kiếm & \texttt{engine/indexer/} \\
M2 & Đánh giá, đo độ trễ, biểu đồ & \texttt{evaluator/}, \texttt{benchmarks/} \\
M3 & Giao diện Streamlit & \texttt{app/} \\
\hline
\end{tabular}
\caption{Các lớp chức năng và ranh giới tích hợp của M2.}
\label{tab:kien-truc-m2}
\end{table}

Phiên bản tích hợp nhận BiMM và bộ lọc tầng 1 từ M1 qua hai callable độc lập.
Baseline âm tiết và fixture tầng 1 chỉ còn dùng cho kiểm thử/phát triển,
được đánh dấu \texttt{development\_only}. Khi nạp chỉ mục,
đối tượng tiền xử lý phải có cùng fingerprint của mã nguồn, cấu hình và tài
nguyên như lúc lập chỉ mục. Cờ tiền xử lý chính thức xác nhận triển khai đã tích
hợp, không thay cho phép đánh giá chất lượng trên gold segmentation do người gán.

\input{chapter3_nlp}

\section{Độ đo TF-IDF và chuẩn hóa}

Với tài liệu $d$ và từ $t$:

\begin{align}
TF(t,d) &= 1+\ln f_{t,d} \quad (f_{t,d}>0), \\
IDF(t) &= \ln\frac{N+1}{df(t)+1}+1, \\
w_{t,d} &= TF(t,d)\times IDF(t), \\
\vec{x}_d &= \frac{w_d}{\lVert w_d\rVert_2}.
\end{align}

Ba lựa chọn cần nêu rõ vì chúng ảnh hưởng trực tiếp tới kết quả đo. Thứ nhất,
dùng sublinear TF nghĩa là một từ xuất hiện ba lần không được coi là gấp ba lần
một từ xuất hiện một lần. Thứ hai, IDF làm mượt $\ln\frac{N+1}{df+1}+1$ không bao
giờ bằng không, nên không mất thông tin ở từ xuất hiện trong mọi tài liệu. Thứ
ba, chuẩn L2 biến tích số vào thành cosin, cho phép so sánh vector có độ dài
khác nhau.

\section{Độ đo TDV và lọc stopword tầng 2}

\subsection{Vì sao cần TDV}

Bộ stopword tĩnh C2 là danh sách tiếng Việt có nguồn và phiên bản cố định. Với
tin tức tiếng Việt theo mười chủ đề, có những từ rất phổ biến nhưng lại chỉ xuất
hiện trong một vài chủ đề; loại chúng đi có thể làm mất thông tin phân biệt. TDV
(Salton) đo trực tiếp câu hỏi đó: bỏ một từ khỏi không gian thì các tài liệu có
dễ phân biệt nhau hơn hay không?

\subsection{Định nghĩa mật độ}

Với $N$ tài liệu đã chuẩn hóa L2, mật độ không gian tài liệu là mật độ cosin
pairwise, loại trừ self-pair:

\begin{equation}
D(X)=\frac{1}{N(N-1)}\sum_{i\neq j}\cos(x_i,x_j).
\end{equation}

Mẫu số $N(N-1)$ là số cặp có thứ tự, được giữ cố định bất kể có bao nhiêu tài
liệu thực sự khác nhau. Với $N<2$ ta đặt $D=0$. Vectơ tất cả bằng không thì
$\cos=0$ với mọi tài liệu khác, nên nó đóng góp bằng không cho mọi cặp chứ không
bị loại khỏi mẫu số.

\subsection{Định nghĩa TDV}

\begin{equation}
TDV(t)=D_{-t}-D_{base},
\end{equation}

trong đó $D_{-t}$ là mật độ sau khi xoá hẳn chiều $t$. Điểm dễ sai nhất nằm
đúng ở chỗ này. Xoá một chiều làm thay đổi hai thứ cùng lúc: độ dài vector của
từng tài liệu, và do đó toàn bộ tổng vector. Vì vậy sau khi xoá chiều, mỗi tài
liệu phải được chuẩn hóa L2 lại, và tổng vector phải được tính lại trên tập
vector đã chuẩn hóa mới.

\subsection{Trọng số đo TDV và tính một lần duy nhất}

TDV được đo trên TF-IDF sơ bộ, tức là trước khi áp bất kỳ bộ lọc TDV nào, và
\textbf{giữ nguyên IDF sơ bộ} khi loại một chiều. Nếu tính lại IDF sau mỗi lần
loại thì hai hiệu ứng trộn vào nhau và con số thu được không còn là TDV của riêng
từ đang xét.

Điểm thứ hai cần nhấn mạnh là TDV được tính \textbf{một lần duy nhất trên toàn
bộ corpus}: mọi từ cùng chia sẻ một ma trận gốc. Không lấy điểm của từ đã bị loại
ở bước trước rồi mới xoá thêm từ sau, vì khi đó điểm đo phụ thuộc thứ tự xoá và
không còn là thuộc tính của từ.

\subsection{Cách cài đặt và kiểm chứng}

Gọi $v_i$ là thành phần chiều $t$ của tài liệu $i$, $S=X X^{\mathsf T}$. Sau khi
xoá chiều $t$, vector của tài liệu $i$ là $x_i - v_i e_t$ với độ dài
$n_i=\sqrt{\lVert x_i\rVert^2-v_i^2}$. Đặt $u_i=1/n_i$, và $u_i=0$ khi vectơ sau
khi xoá bằng không. Khi đó tổng cosin ngoài đường chéo sau khi xoá là

\begin{equation}
\sum_{i\neq j} u_iu_j\,(S_{ij}-v_iv_j)
= \lVert X^{\mathsf T}u\rVert^2 - \sum_i u_i^2 S_{ii}
- \Bigl(\langle u,v\rangle^2-\sum_i u_i^2 v_i^2\Bigr).
\end{equation}

Nhờ đẳng thức này, mỗi từ chỉ cần duyệt posting list của nó thay vì dựng lại ma
trận cosin đầy đủ $N\times N$.

Cách làm sai mà đề cương cảnh báo, tức là xoá chiều mà không tái chuẩn hóa và
không cập nhật tổng vector, được cài sẵn trong bộ kiểm thử dưới dạng một hàm tham
chiếu thứ hai. Bộ test khẳng định kết quả của engine khác kết quả của cách làm
sai, đồng thời khớp với một hàm tham chiếu thứ ba hoàn toàn trực tiếp: xoá cột,
tái chuẩn hóa từng dòng, nhân ma trận cosin đầy đủ.

\subsection{Ngưỡng lọc}

\begin{equation}
\text{giữ từ } t \iff TDV(t)\geq -\varepsilon,\qquad \varepsilon=10^{-12}.
\end{equation}

$\varepsilon$ là dung sai số học chứ không phải siêu tham số thuật toán truy xuất;
giá trị này được ghi vào metadata của mọi lần chạy. Không ép đường cong hình chữ
U, và không suy ra từ dấu của TDV rằng một cấu hình lọc nào sẽ cho kết quả tốt
hơn: việc loại một từ có cải thiện truy xuất hay không chỉ trả lời được bằng phép
đo với qrels đã gán nhãn.

\section{Bốn cấu hình ablation}

\begin{center}
\begin{tabular}{|l|l|l|}
\hline
Mã & Cấu hình & Ghi chú \\
\hline
C1 & Không lọc stopword & baseline \\
C2 & Stoplist tĩnh, độc lập với corpus & nguồn và số mục thực tế được ghi kèm \\
C3 & Chỉ stopword ngữ pháp tầng 1 của M1 & \\
C4 & Tầng 1 cộng TDV tầng 2 & giữ các từ phủ định bằng guard ngữ nghĩa \\
\hline
\end{tabular}
\end{center}

Mỗi cấu hình nằm trong một thư mục riêng kèm metadata: SHA256 corpus, tokenizer,
cấu hình, stoplist và số phần tử stoplist. Bốn chỉ mục dùng chung một corpus, một
bộ truy vấn và một bộ nhãn, nên so sánh giữa chúng là so sánh đúng yếu tố lọc
stopword.

\subsection{Về lớp tầng 1 tạm thời}

Hợp đồng dự án quy định: vì đồ án chưa có Maximum Matching chính thức của M1, nếu
C3/C4 dùng bộ lọc tầng 1 tạm thì kết quả phải được ghi là provisional, không
coi là Maximum Matching chính thức. Bộ đánh giá vì vậy chặn tuyệt đối mọi kết quả
chính thức khi metadata còn trỏ tới bộ lọc tạm hoặc tokenizer phát triển, kể cả
khi metadata ghi nhầm trạng thái là official.

\subsection{Giữ từ phủ định}

Bộ TDV theo bản chất thiên về việc loại từ phổ biến, mà các từ phủ định như
``không'', ``chưa'', ``chẳng'', ``vì'' lại mang thông tin quan trọng cho truy vấn
dạng phủ định. Vì vậy sau khi TDV chọn stoplist, một danh sách guard ngữ nghĩa
được gỡ lại khỏi stoplist. Với corpus chính, TDV chọn 498 từ và sau bước guard còn
494 từ, với bốn từ được giữ lại là ``chưa'', ``chẳng'', ``không'' và ``vì''.

\subsection{Cơ chế fallback của C4 và giới hạn của nó}

Khi một truy vấn mà mọi token đều đã bị lọc khỏi từ vựng của C4, vectơ truy vấn
bằng không và không thể có kết quả. Cơ chế dự phòng là tra cứu bằng một chỉ mục
C3 riêng, có cùng corpus và cùng tiền xử lý. Cần nói rõ hai điều. Một, giữ lại một
từ đã bị xoá khỏi chỉ mục không phục hồi được khả năng truy xuất bằng chỉ mục đó;
fallback chỉ cứu được truy vấn mà không còn token nào. Hai, hệ thống khi đó phải
giữ cả hai chỉ mục, nên dung lượng triển khai tăng lên chứ không giảm; số liệu
mục 6.1 của Chương 4 đo riêng cả chỉ mục C4 thuần lẫn gói C4 cộng C3.

\section{Cấu trúc chỉ mục và đường đi của một truy vấn}

\begin{enumerate}
\item Chuẩn hóa văn bản và tách từ theo đối tượng tiền xử lý của M1.
\item Lọc stopword theo đúng cấu hình đã dùng khi lập chỉ mục.
\item Vector hoá truy vấn bằng bảng IDF ngoại tuyến của corpus.
\item Lấy ứng viên bằng hợp các posting list trong chỉ mục đảo.
\item Tính cosin trên tập ứng viên, chấm điểm chỉ những kết quả có điểm dương.
\item Sắp xếp giảm dần điểm; hòa điểm sắp theo \texttt{doc\_id} tăng dần để kết
quả tái lập được.
\end{enumerate}

Chỉ mục được lưu dạng thưa: posting list và vector trọng số. Hàm giải thích
(\texttt{explain}) trả về từng đóng góp vào tích vô hướng, và tổng các đóng góp
bằng đúng điểm cosin đã trả về, để người đọc có thể kiểm tra lại cách xếp hạng.

\section{Giao diện đánh giá và cổng chặn kết quả chính thức}

Bộ đánh giá dùng các định nghĩa chỉ số sau, với $K$ mặc định là 10 và có thể cấu
hình:

\begin{align}
P@K &= \frac{\text{số trúng trong Top-}K}{K}, \\
Recall@K &= \frac{\text{số trúng trong Top-}K}{\text{tổng số relevant đã confirmed}}, \\
F1@K &= \frac{2 \cdot P@K \cdot Recall@K}{P@K + Recall@K}, \\
AP &= \frac{1}{R}\sum_{i=1}^{L} P@i \cdot rel_i, \\
NDCG@K &= \frac{DCG}{IDCG},\quad DCG=\sum_{i=1}^{K}\frac{2^{rel_i}-1}{\log_2(i+1)} .
\end{align}

Ba quyết định thiết kế đáng ghi lại. Mẫu số của $P@K$ là $K$ cố định, kể cả khi
hệ thống trả về ít hơn $K$ kết quả; thiếu kết quả là sai chứ không phải là không
tính. $AP$ tính trên toàn bộ ranking có sẵn, không cắt theo $K$, nên bộ đánh giá
từ chối chạy nếu độ sâu lấy ranking nhỏ hơn số tài liệu corpus. $IDCG$ lấy từ
qrels chứ không lấy từ lần chạy, nên điểm không thể bị đẩy lên bằng cách chỉ
trả về một phần kết quả.

Query không có tài liệu relevant nào vẫn được giữ trong bình quân với giá trị 0 và
được đếm riêng, để người đọc biết chính xác MAP bị kéo xuống bao nhiêu do những
query đó. Tương tự, một truy vấn không có kết quả trả về vẫn sinh ra dòng chỉ số
bằng 0 chứ không bị loại khỏi bình quân.

\subsection{Điều kiện để coi là đánh giá chính thức}

Bộ đánh gia từ chối ghi kết quả chính thức khi có bất kỳ điều kiện nào sau đây:

\begin{itemize}
\item dưới 20 truy vấn;
\item SHA256 corpus lệch giữa manifest, qrels, tập truy vấn hoặc chỉ mục, hoặc
thiếu hẳn hash ở bất kỳ nguồn nào;
\item còn phán quyết chưa confirmed, hoặc nhãn confirmed mà điểm relevance còn
null, hoặc nhãn confirmed không ghi người gán;
\item trùng cặp (truy vấn, tài liệu), điểm relevance ngoài tập $\{0,1,2\}$, hoặc
truy vấn trong qrels không tồn tại trong tập truy vấn;
\item metadata tiền xử lý chưa nhất quán ở trạng thái official, hoặc đang dùng
tokenizer phát triển hay bộ lọc tầng 1 tạm.
\end{itemize}

SHA256 corpus cần đánh giá không bao giờ được lấy từ qrels: nếu không xác định
được từ manifest đã khóa hoặc từ chính file corpus thì kết quả chính thức bị chặn.

\subsection{Khi chưa đủ nhãn thì không xuất chỉ số}

Nếu bất kỳ truy vấn nào chưa có nhãn confirmed, bộ đánh giá ghi
\texttt{metrics\_status = unmeasured} và không ghi bảng chỉ số. Lý do rất cụ thể:
với qrels rỗng, mọi chỉ số đều bằng không một cách máy móc, và một bảng toàn số
không dễ bị đọc nhầm là hệ thống không tìm được gì.

\section{Tách kết quả kiểm thử khỏi kết quả chính thức}

Các bài kiểm thử dùng dữ liệu tạm trong thư mục làm việc của pytest. Để ngăn
tình huống số liệu kiểm thử lọt vào báo cáo, bộ đánh giá chặn ghi kết quả không
chính thức vào các thư mục kết quả chính thức; kết quả chạy thử bắt buộc ghi ra
thư mục khác, ví dụ \texttt{results/development\_eval}.
 trong main.tex.
% Chỉ dùng các gói đã khai báo trong main.tex: report, vietnamese, amsmath,
% amssymb, graphicx, hyperref, geometry. Không thêm gói mới.
\chapter{THIẾT KẾ HỆ THỐNG TÌM KIẾM VĂN BẢN}

\section{Phạm vi trách nhiệm của M2 trong kiến trúc}

Kiến trúc tổng thể chia thành ba lớp trách nhiệm. M1 sở hữu tiền xử lý ngôn ngữ
tự nhiên, M2 sở hữu biến các từ đã tách thành dữ liệu lập chỉ mục và truy xuất,
M3 sở hữu giao diện người dùng. Chương này trình bày phần của M2.

Các lớp và ranh giới tích hợp được trình bày trong Bảng~\ref{tab:kien-truc-m2}.
\begin{table}[ht]
\centering
\begin{tabular}{|p{0.08\textwidth}|p{0.43\textwidth}|p{0.35\textwidth}|}
\hline
Lớp & Thành phần & File chính \\
\hline
M1 & Maximum Matching, stopword tầng 1 & \texttt{engine/tokenizer/} \\
M2 & Adapter M1 có fingerprint & \texttt{engine/preprocessing.py} \\
M2 & TDV, TF-IDF, chỉ mục đảo, tìm kiếm & \texttt{engine/indexer/} \\
M2 & Đánh giá, đo độ trễ, biểu đồ & \texttt{evaluator/}, \texttt{benchmarks/} \\
M3 & Giao diện Streamlit & \texttt{app/} \\
\hline
\end{tabular}
\caption{Các lớp chức năng và ranh giới tích hợp của M2.}
\label{tab:kien-truc-m2}
\end{table}

Phiên bản tích hợp nhận BiMM và bộ lọc tầng 1 từ M1 qua hai callable độc lập.
Baseline âm tiết và fixture tầng 1 chỉ còn dùng cho kiểm thử/phát triển,
được đánh dấu \texttt{development\_only}. Khi nạp chỉ mục,
đối tượng tiền xử lý phải có cùng fingerprint của mã nguồn, cấu hình và tài
nguyên như lúc lập chỉ mục. Cờ tiền xử lý chính thức xác nhận triển khai đã tích
hợp, không thay cho phép đánh giá chất lượng trên gold segmentation do người gán.

% Phần M1 cho Chương 3 — được \input bên trong chapter3_ir.tex nên KHÔNG có \chapter.
% Chỉ dùng các gói đã khai báo trong main.tex; không thêm gói mới.
\section{Thiết kế tầng ngôn ngữ M1}
\label{sec:m1-thiet-ke}

Tầng ngôn ngữ M1 biến văn bản thô thành dãy token sạch để M2 lập chỉ mục. Tầng này
gồm ba khối độc lập: chuẩn hóa chuỗi, tách từ theo so khớp cực đại trên từ điển, và
lọc từ dừng tầng~1 kèm danh sách bảo vệ. Toàn bộ tầng hoạt động ngoại tuyến, không
gọi dịch vụ ngoài, và mọi tài nguyên mà nó dùng đều được ghim bằng hàm băm.

\subsection{Chuẩn hóa văn bản và bảo toàn ranh giới}

Hàm chuẩn hóa \texttt{normalize(text)} trả về dạng hiển thị: chuẩn hóa Unicode
NFC, hạ chữ thường, loại bỏ thẻ đánh dấu, đường dẫn và địa chỉ thư điện tử, giữ
lại chữ cái tiếng Việt cùng khoảng trắng, và thu gọn mọi khoảng trắng về một dấu
cách. Một ví dụ:

\begin{quote}
\texttt{"Tri thức, nhân loại! 2026"} $\rightarrow$ \texttt{"tri thức nhân loại"}.
\end{quote}

Một quyết định thiết kế quan trọng là \emph{dấu câu là ranh giới, không phải
nhiễu}. Dạng hiển thị gộp khoảng trắng nên trông như một chuỗi liên tục, nhưng bộ
tách từ phải không bao giờ ghép một từ xuyên qua dấu câu. Vì lý do đó, bộ tách từ
tiêu thụ một dạng trung gian giữ lại ranh giới, ở đó đoạn văn bản bị cắt tại mỗi
dấu câu và tại mỗi khoảng trắng không phải dấu cách thường (xuống dòng, tab). Các
khoảng được gỡ bỏ (đường dẫn, thư điện tử, thẻ đánh dấu) cũng được thay bằng một
ranh giới cứng chứ không phải bằng dấu cách; nếu thay bằng dấu cách thì hai mảnh
văn bản hai bên sẽ bị dính thành một từ, ví dụ \texttt{"học<br>sinh"} phải tách
thành hai âm tiết chứ không được ghép thành một từ.

Cũng trong khối này có một chỉnh sửa so với mẫu biểu thức chính quy ban đầu. Mẫu
cũ chỉ giữ chữ cái tiếng Việt nên đã vô tình xóa mọi chữ cái Latin, biến
``Windows'' thành ``indows'' và ``Java'' thành ``ava''. Bốn chữ cái Latin
\texttt{f}, \texttt{j}, \texttt{w}, \texttt{z} được giữ lại vì chúng phân biệt
trong các tên vay mượn phổ biến (Windows, WWW, Java, Zalo). Chúng \emph{không}
được thêm vào từ điển; kết quả trung thực là chúng trở thành token đơn ngoài từ
điển thay vì bị bóp méo thành từ sai.

\subsection{Từ điển và cây Trie}

Tách từ dựa trên từ điển Viet74K \cite{viet74k} gồm 69.858 mục, được ghim hàm băm
SHA256 \texttt{175e8b1b...d8ab} và đi kèm thông tin giấy phép cùng nguồn gốc trong
kho. Từ điển được nạp vào một cây Trie hỗ trợ truy vấn theo tiền tố, cho phép bộ
tách từ kiểm tra đoạn ứng viên có trong từ điển hay không mà không phải so khớp
toàn bộ danh sách. Cửa sổ trượt mặc định lấy bằng độ dài mục dài nhất trong từ
điển, nên không bỏ sót từ nào; người dùng vẫn có thể chỉ định một giới hạn cửa sổ
nhỏ hơn khi cần quan sát thuật toán.

\subsection{Thuật toán FMM, BMM và BiMM}

Ba phương pháp dùng chung một cửa sổ trượt nhưng khác hướng duyệt. FMM duyệt tiến,
BMM duyệt lùi. BiMM chạy cả hai rồi chọn theo bộ quy tắc phân xử sau, xét theo
thứ tự:

\begin{enumerate}
  \item ít token hơn;
  \item ít token một âm tiết hơn;
  \item ít token ngoài từ điển hơn;
  \item nếu vẫn hòa, chọn BMM.
\end{enumerate}

Tiêu chí thứ ba là lưới an toàn tường minh: với cài đặt Trie thuần, mọi token
ngoài từ điển đều được phát ra dưới dạng một âm tiết, nên tiêu chí này được suy ra
từ tiêu chí thứ hai; nhóm vẫn giữ nó cùng khóa so sánh đầy đủ để người đọc kiểm
chứng lại quyết định thay vì tin vào phần tóm tắt.

Ví dụ nhập nhằng kinh điển minh họa vì sao cần phân xử. Với câu ``học sinh học
sinh học'' và một từ điển fixture có cả ``học sinh'' lẫn ``sinh học'':

\begin{quote}
FMM: \texttt{học\_sinh, học\_sinh, học}\\
BMM: \texttt{học, sinh\_học, sinh\_học}\\
BiMM: \texttt{học, sinh\_học, sinh\_học}
\end{quote}

Hai phương án có cùng số token và cùng số token một âm tiết, nên hòa và BiMM chọn
BMM. Ví dụ này dùng một từ điển fixture minh bạch, không phụ thuộc từ điển đóng
gói, để phép so sánh là một khẳng định về \emph{thuật toán} chứ không về một từ
điển cụ thể.

Mỗi lần chạy đều sinh một \emph{vết} từng bước. Mỗi dòng vết ghi số thứ tự bước,
hướng duyệt, chỉ số âm tiết bắt đầu và kết thúc (kết thúc mở), đoạn ứng viên, độ
dài cửa sổ, cờ có khớp từ điển hay không, và hành động -- một trong
\texttt{MATCHED}, \texttt{SHRINK} hoặc \texttt{SINGLE}. Vết dùng đúng từ điển và
cửa sổ của lần chạy, nên người đọc có thể tái lập từng quyết định thay vì chỉ thấy
kết quả cuối cùng. Với đầu vào rỗng hoặc chỉ gồm dấu câu, kết quả là dãy token rỗng
và vết rỗng.

\subsection{Lọc từ dừng tầng 1 và danh sách bảo vệ}

Khối thứ ba loại các hư từ ngữ pháp -- đại từ, giới từ, trợ từ, phó từ chỉ mức độ,
liên từ -- vốn đóng góp ít cho phân biệt truy xuất. Hai tài nguyên có nguồn và quy
tắc biên soạn điều khiển bước này: danh sách hư từ ngữ pháp và danh sách bảo vệ
(guard). Tập từ dừng hiệu lực là hiệu của hai danh sách: \emph{guard luôn thắng}.

Quy tắc bảo vệ có phạm vi rộng hơn một so khớp đơn giản: một token được bảo vệ nếu
chính nó là guard, hoặc nếu một thành phần của nó (tách theo dấu gạch nối âm tiết)
cũng là guard. Nhờ đó việc bảo vệ ``không'' cũng giữ nguyên ``không\_thể'', bảo vệ
``vì'' giữ nguyên ``vì\_dầu'', bảo vệ ``chẳng'' giữ nguyên ``chẳng\_hạn''. Điểm
tinh tế là chỉ những thành phần \emph{tự chúng là guard} mới có thể bảo vệ: nếu mọi
âm tiết của mọi guard đều có thể bảo vệ thì ``thể'' sẽ được bảo vệ và ``thể\_thao''
bị bảo vệ nhầm. Chính vì vậy các từ ghép cần bảo vệ vì lý do khác được liệt kê
tường minh trong tệp guard.

Ví dụ minh họa:

\begin{quote}
\texttt{filter\_tier1(["bóng", "đá", "và", "của", "bóng\_đá"])}
$\rightarrow$ \texttt{["bóng", "đá", "bóng\_đá"]}
\end{quote}

Ở đây ``và'' và ``của'' bị loại, còn ``bóng\_đá'' là một từ ghép nội dung được giữ
nguyên. Cùng cơ chế bảo vệ này được tái dùng cho tầng~2 khi TDV chọn stoplist, để
các từ phủ định không bị xóa ở cả hai tầng.

\subsection{Tài nguyên và xuất xứ}

Ba tài nguyên ngôn ngữ học được đóng gói kèm số mục thực tế và hàm băm, để kết quả
tái lập được và để bất kỳ ai cũng kiểm tra lại được dữ liệu đầu vào.

\begin{table}[ht]
\centering
\begin{tabular}{|p{0.30\textwidth}|p{0.38\textwidth}|p{0.18\textwidth}|}
\hline
Tài nguyên & Vai trò & Số mục \\
\hline
\texttt{dict.txt} & Từ điển tách từ (Viet74K) & 69.858 \\
\texttt{grammar\_\allowbreak stopwords.txt} & Hư từ ngữ pháp tầng 1 & 181 \\
\texttt{guard\_\allowbreak whitelist.txt} & Danh sách bảo vệ (guard) & 110 \\
\hline
\end{tabular}
\caption{Ba tài nguyên của tầng ngôn ngữ M1 (đặt trong thư mục \texttt{resources/})
và số mục tại thời điểm viết. Số mục được đọc trực tiếp từ tệp; danh sách hư từ sau
khi trừ guard còn 124 từ dừng hiệu lực.}
\label{tab:m1-tai-nguyen}
\end{table}

Từ điển là danh sách mã nguồn mở có ghi giấy phép \cite{viet74k}. Danh sách từ
dừng tĩnh dùng cho cấu hình C2 là một bản chụp có nguồn, phiên bản và thông báo
giấy phép MIT \cite{static_stopwords}. Hai danh sách tầng~1 do nhóm biên soạn kèm
quy tắc và số mục thực tế; các từ dừng được gỡ lại chỉ là các từ phủ định và logic
có lý do ngữ nghĩa rõ ràng, không bù danh sách bằng thực từ để đạt một con số.

\subsection{Kiểm chứng tầng ngôn ngữ}

Ba phương pháp tách từ được kiểm chứng bằng từ điển fixture minh bạch và bằng một
đường tham chiếu trực tiếp trên từ điển/câu ngẫu nhiên. Vết thuật toán cho biết ứng
viên và quyết định tại từng bước; dãy token do \texttt{tokenize} trả về được kiểm
riêng bằng fixture, vì với BMM các dòng ứng viên được đánh dấu có thể chồng lấn nên
không dùng chúng để tái dựng kết quả. Bộ lọc tầng~1 được kiểm chứng
bằng các ca guard: phủ định và từ ghép chứa guard phải được giữ nguyên, còn hư từ
không nằm trong guard phải bị loại. Các ca kiểm thử này chứng minh \emph{luật hoạt
động đúng}, không đo \emph{chất lượng} tách từ.

Chất lượng tách từ chỉ đo được khi có nhãn người. Nhóm đã dựng bộ mẫu 150 câu
(15 câu mỗi lĩnh vực) rút từ kho ngữ liệu đã khóa, kèm công cụ gán nhãn ngoại
tuyến và bộ đo khớp span từ. Hiện 150/150 câu có cờ xác nhận và phép đo word-span
đã chạy, nhưng ba câu SG010, SG053, SG089 còn cần rà lại về ranh giới từ; do đó
F1 ở Chương~4 là số đo tạm thời trên gold hiện có, không phải kết luận chất lượng
cuối. Gợi ý của máy chỉ phục vụ gán nhãn nhanh và không tự động trở thành gold.
Quy trình đo, các ràng buộc theo hàm băm và giới hạn của phép đo trình bày trong
Chương~4.


\section{Độ đo TF-IDF và chuẩn hóa}

Với tài liệu $d$ và từ $t$:

\begin{align}
TF(t,d) &= 1+\ln f_{t,d} \quad (f_{t,d}>0), \\
IDF(t) &= \ln\frac{N+1}{df(t)+1}+1, \\
w_{t,d} &= TF(t,d)\times IDF(t), \\
\vec{x}_d &= \frac{w_d}{\lVert w_d\rVert_2}.
\end{align}

Ba lựa chọn cần nêu rõ vì chúng ảnh hưởng trực tiếp tới kết quả đo. Thứ nhất,
dùng sublinear TF nghĩa là một từ xuất hiện ba lần không được coi là gấp ba lần
một từ xuất hiện một lần. Thứ hai, IDF làm mượt $\ln\frac{N+1}{df+1}+1$ không bao
giờ bằng không, nên không mất thông tin ở từ xuất hiện trong mọi tài liệu. Thứ
ba, chuẩn L2 biến tích số vào thành cosin, cho phép so sánh vector có độ dài
khác nhau.

\section{Độ đo TDV và lọc stopword tầng 2}

\subsection{Vì sao cần TDV}

Bộ stopword tĩnh C2 là danh sách tiếng Việt có nguồn và phiên bản cố định. Với
tin tức tiếng Việt theo mười chủ đề, có những từ rất phổ biến nhưng lại chỉ xuất
hiện trong một vài chủ đề; loại chúng đi có thể làm mất thông tin phân biệt. TDV
(Salton) đo trực tiếp câu hỏi đó: bỏ một từ khỏi không gian thì các tài liệu có
dễ phân biệt nhau hơn hay không?

\subsection{Định nghĩa mật độ}

Với $N$ tài liệu đã chuẩn hóa L2, mật độ không gian tài liệu là mật độ cosin
pairwise, loại trừ self-pair:

\begin{equation}
D(X)=\frac{1}{N(N-1)}\sum_{i\neq j}\cos(x_i,x_j).
\end{equation}

Mẫu số $N(N-1)$ là số cặp có thứ tự, được giữ cố định bất kể có bao nhiêu tài
liệu thực sự khác nhau. Với $N<2$ ta đặt $D=0$. Vectơ tất cả bằng không thì
$\cos=0$ với mọi tài liệu khác, nên nó đóng góp bằng không cho mọi cặp chứ không
bị loại khỏi mẫu số.

\subsection{Định nghĩa TDV}

\begin{equation}
TDV(t)=D_{-t}-D_{base},
\end{equation}

trong đó $D_{-t}$ là mật độ sau khi xoá hẳn chiều $t$. Điểm dễ sai nhất nằm
đúng ở chỗ này. Xoá một chiều làm thay đổi hai thứ cùng lúc: độ dài vector của
từng tài liệu, và do đó toàn bộ tổng vector. Vì vậy sau khi xoá chiều, mỗi tài
liệu phải được chuẩn hóa L2 lại, và tổng vector phải được tính lại trên tập
vector đã chuẩn hóa mới.

\subsection{Trọng số đo TDV và tính một lần duy nhất}

TDV được đo trên TF-IDF sơ bộ, tức là trước khi áp bất kỳ bộ lọc TDV nào, và
\textbf{giữ nguyên IDF sơ bộ} khi loại một chiều. Nếu tính lại IDF sau mỗi lần
loại thì hai hiệu ứng trộn vào nhau và con số thu được không còn là TDV của riêng
từ đang xét.

Điểm thứ hai cần nhấn mạnh là TDV được tính \textbf{một lần duy nhất trên toàn
bộ corpus}: mọi từ cùng chia sẻ một ma trận gốc. Không lấy điểm của từ đã bị loại
ở bước trước rồi mới xoá thêm từ sau, vì khi đó điểm đo phụ thuộc thứ tự xoá và
không còn là thuộc tính của từ.

\subsection{Cách cài đặt và kiểm chứng}

Gọi $v_i$ là thành phần chiều $t$ của tài liệu $i$, $S=X X^{\mathsf T}$. Sau khi
xoá chiều $t$, vector của tài liệu $i$ là $x_i - v_i e_t$ với độ dài
$n_i=\sqrt{\lVert x_i\rVert^2-v_i^2}$. Đặt $u_i=1/n_i$, và $u_i=0$ khi vectơ sau
khi xoá bằng không. Khi đó tổng cosin ngoài đường chéo sau khi xoá là

\begin{equation}
\sum_{i\neq j} u_iu_j\,(S_{ij}-v_iv_j)
= \lVert X^{\mathsf T}u\rVert^2 - \sum_i u_i^2 S_{ii}
- \Bigl(\langle u,v\rangle^2-\sum_i u_i^2 v_i^2\Bigr).
\end{equation}

Nhờ đẳng thức này, mỗi từ chỉ cần duyệt posting list của nó thay vì dựng lại ma
trận cosin đầy đủ $N\times N$.

Cách làm sai mà đề cương cảnh báo, tức là xoá chiều mà không tái chuẩn hóa và
không cập nhật tổng vector, được cài sẵn trong bộ kiểm thử dưới dạng một hàm tham
chiếu thứ hai. Bộ test khẳng định kết quả của engine khác kết quả của cách làm
sai, đồng thời khớp với một hàm tham chiếu thứ ba hoàn toàn trực tiếp: xoá cột,
tái chuẩn hóa từng dòng, nhân ma trận cosin đầy đủ.

\subsection{Ngưỡng lọc}

\begin{equation}
\text{giữ từ } t \iff TDV(t)\geq -\varepsilon,\qquad \varepsilon=10^{-12}.
\end{equation}

$\varepsilon$ là dung sai số học chứ không phải siêu tham số thuật toán truy xuất;
giá trị này được ghi vào metadata của mọi lần chạy. Không ép đường cong hình chữ
U, và không suy ra từ dấu của TDV rằng một cấu hình lọc nào sẽ cho kết quả tốt
hơn: việc loại một từ có cải thiện truy xuất hay không chỉ trả lời được bằng phép
đo với qrels đã gán nhãn.

\section{Bốn cấu hình ablation}

\begin{center}
\begin{tabular}{|l|l|l|}
\hline
Mã & Cấu hình & Ghi chú \\
\hline
C1 & Không lọc stopword & baseline \\
C2 & Stoplist tĩnh, độc lập với corpus & nguồn và số mục thực tế được ghi kèm \\
C3 & Chỉ stopword ngữ pháp tầng 1 của M1 & \\
C4 & Tầng 1 cộng TDV tầng 2 & giữ các từ phủ định bằng guard ngữ nghĩa \\
\hline
\end{tabular}
\end{center}

Mỗi cấu hình nằm trong một thư mục riêng kèm metadata: SHA256 corpus, tokenizer,
cấu hình, stoplist và số phần tử stoplist. Bốn chỉ mục dùng chung một corpus, một
bộ truy vấn và một bộ nhãn, nên so sánh giữa chúng là so sánh đúng yếu tố lọc
stopword.

\subsection{Về lớp tầng 1 tạm thời}

Hợp đồng dự án quy định: vì đồ án chưa có Maximum Matching chính thức của M1, nếu
C3/C4 dùng bộ lọc tầng 1 tạm thì kết quả phải được ghi là provisional, không
coi là Maximum Matching chính thức. Bộ đánh giá vì vậy chặn tuyệt đối mọi kết quả
chính thức khi metadata còn trỏ tới bộ lọc tạm hoặc tokenizer phát triển, kể cả
khi metadata ghi nhầm trạng thái là official.

\subsection{Giữ từ phủ định}

Bộ TDV theo bản chất thiên về việc loại từ phổ biến, mà các từ phủ định như
``không'', ``chưa'', ``chẳng'', ``vì'' lại mang thông tin quan trọng cho truy vấn
dạng phủ định. Vì vậy sau khi TDV chọn stoplist, một danh sách guard ngữ nghĩa
được gỡ lại khỏi stoplist. Với corpus chính, TDV chọn 498 từ và sau bước guard còn
494 từ, với bốn từ được giữ lại là ``chưa'', ``chẳng'', ``không'' và ``vì''.

\subsection{Cơ chế fallback của C4 và giới hạn của nó}

Khi một truy vấn mà mọi token đều đã bị lọc khỏi từ vựng của C4, vectơ truy vấn
bằng không và không thể có kết quả. Cơ chế dự phòng là tra cứu bằng một chỉ mục
C3 riêng, có cùng corpus và cùng tiền xử lý. Cần nói rõ hai điều. Một, giữ lại một
từ đã bị xoá khỏi chỉ mục không phục hồi được khả năng truy xuất bằng chỉ mục đó;
fallback chỉ cứu được truy vấn mà không còn token nào. Hai, hệ thống khi đó phải
giữ cả hai chỉ mục, nên dung lượng triển khai tăng lên chứ không giảm; số liệu
mục 6.1 của Chương 4 đo riêng cả chỉ mục C4 thuần lẫn gói C4 cộng C3.

\section{Cấu trúc chỉ mục và đường đi của một truy vấn}

\begin{enumerate}
\item Chuẩn hóa văn bản và tách từ theo đối tượng tiền xử lý của M1.
\item Lọc stopword theo đúng cấu hình đã dùng khi lập chỉ mục.
\item Vector hoá truy vấn bằng bảng IDF ngoại tuyến của corpus.
\item Lấy ứng viên bằng hợp các posting list trong chỉ mục đảo.
\item Tính cosin trên tập ứng viên, chấm điểm chỉ những kết quả có điểm dương.
\item Sắp xếp giảm dần điểm; hòa điểm sắp theo \texttt{doc\_id} tăng dần để kết
quả tái lập được.
\end{enumerate}

Chỉ mục được lưu dạng thưa: posting list và vector trọng số. Hàm giải thích
(\texttt{explain}) trả về từng đóng góp vào tích vô hướng, và tổng các đóng góp
bằng đúng điểm cosin đã trả về, để người đọc có thể kiểm tra lại cách xếp hạng.

\section{Giao diện đánh giá và cổng chặn kết quả chính thức}

Bộ đánh giá dùng các định nghĩa chỉ số sau, với $K$ mặc định là 10 và có thể cấu
hình:

\begin{align}
P@K &= \frac{\text{số trúng trong Top-}K}{K}, \\
Recall@K &= \frac{\text{số trúng trong Top-}K}{\text{tổng số relevant đã confirmed}}, \\
F1@K &= \frac{2 \cdot P@K \cdot Recall@K}{P@K + Recall@K}, \\
AP &= \frac{1}{R}\sum_{i=1}^{L} P@i \cdot rel_i, \\
NDCG@K &= \frac{DCG}{IDCG},\quad DCG=\sum_{i=1}^{K}\frac{2^{rel_i}-1}{\log_2(i+1)} .
\end{align}

Ba quyết định thiết kế đáng ghi lại. Mẫu số của $P@K$ là $K$ cố định, kể cả khi
hệ thống trả về ít hơn $K$ kết quả; thiếu kết quả là sai chứ không phải là không
tính. $AP$ tính trên toàn bộ ranking có sẵn, không cắt theo $K$, nên bộ đánh giá
từ chối chạy nếu độ sâu lấy ranking nhỏ hơn số tài liệu corpus. $IDCG$ lấy từ
qrels chứ không lấy từ lần chạy, nên điểm không thể bị đẩy lên bằng cách chỉ
trả về một phần kết quả.

Query không có tài liệu relevant nào vẫn được giữ trong bình quân với giá trị 0 và
được đếm riêng, để người đọc biết chính xác MAP bị kéo xuống bao nhiêu do những
query đó. Tương tự, một truy vấn không có kết quả trả về vẫn sinh ra dòng chỉ số
bằng 0 chứ không bị loại khỏi bình quân.

\subsection{Điều kiện để coi là đánh giá chính thức}

Bộ đánh gia từ chối ghi kết quả chính thức khi có bất kỳ điều kiện nào sau đây:

\begin{itemize}
\item dưới 20 truy vấn;
\item SHA256 corpus lệch giữa manifest, qrels, tập truy vấn hoặc chỉ mục, hoặc
thiếu hẳn hash ở bất kỳ nguồn nào;
\item còn phán quyết chưa confirmed, hoặc nhãn confirmed mà điểm relevance còn
null, hoặc nhãn confirmed không ghi người gán;
\item trùng cặp (truy vấn, tài liệu), điểm relevance ngoài tập $\{0,1,2\}$, hoặc
truy vấn trong qrels không tồn tại trong tập truy vấn;
\item metadata tiền xử lý chưa nhất quán ở trạng thái official, hoặc đang dùng
tokenizer phát triển hay bộ lọc tầng 1 tạm.
\end{itemize}

SHA256 corpus cần đánh giá không bao giờ được lấy từ qrels: nếu không xác định
được từ manifest đã khóa hoặc từ chính file corpus thì kết quả chính thức bị chặn.

\subsection{Khi chưa đủ nhãn thì không xuất chỉ số}

Nếu bất kỳ truy vấn nào chưa có nhãn confirmed, bộ đánh giá ghi
\texttt{metrics\_status = unmeasured} và không ghi bảng chỉ số. Lý do rất cụ thể:
với qrels rỗng, mọi chỉ số đều bằng không một cách máy móc, và một bảng toàn số
không dễ bị đọc nhầm là hệ thống không tìm được gì.

\section{Tách kết quả kiểm thử khỏi kết quả chính thức}

Các bài kiểm thử dùng dữ liệu tạm trong thư mục làm việc của pytest. Để ngăn
tình huống số liệu kiểm thử lọt vào báo cáo, bộ đánh giá chặn ghi kết quả không
chính thức vào các thư mục kết quả chính thức; kết quả chạy thử bắt buộc ghi ra
thư mục khác, ví dụ \texttt{results/development\_eval}.
 trong main.tex.
% Chỉ dùng các gói đã khai báo trong main.tex: report, vietnamese, amsmath,
% amssymb, graphicx, hyperref, geometry. Không thêm gói mới.
\chapter{THIẾT KẾ HỆ THỐNG TÌM KIẾM VĂN BẢN}

\section{Phạm vi trách nhiệm của M2 trong kiến trúc}

Kiến trúc tổng thể chia thành ba lớp trách nhiệm. M1 sở hữu tiền xử lý ngôn ngữ
tự nhiên, M2 sở hữu biến các từ đã tách thành dữ liệu lập chỉ mục và truy xuất,
M3 sở hữu giao diện người dùng. Chương này trình bày phần của M2.

Các lớp và ranh giới tích hợp được trình bày trong Bảng~\ref{tab:kien-truc-m2}.
\begin{table}[ht]
\centering
\begin{tabular}{|p{0.08\textwidth}|p{0.43\textwidth}|p{0.35\textwidth}|}
\hline
Lớp & Thành phần & File chính \\
\hline
M1 & Maximum Matching, stopword tầng 1 & \texttt{engine/tokenizer/} \\
M2 & Adapter M1 có fingerprint & \texttt{engine/preprocessing.py} \\
M2 & TDV, TF-IDF, chỉ mục đảo, tìm kiếm & \texttt{engine/indexer/} \\
M2 & Đánh giá, đo độ trễ, biểu đồ & \texttt{evaluator/}, \texttt{benchmarks/} \\
M3 & Giao diện Streamlit & \texttt{app/} \\
\hline
\end{tabular}
\caption{Các lớp chức năng và ranh giới tích hợp của M2.}
\label{tab:kien-truc-m2}
\end{table}

Phiên bản tích hợp nhận BiMM và bộ lọc tầng 1 từ M1 qua hai callable độc lập.
Baseline âm tiết và fixture tầng 1 chỉ còn dùng cho kiểm thử/phát triển,
được đánh dấu \texttt{development\_only}. Khi nạp chỉ mục,
đối tượng tiền xử lý phải có cùng fingerprint của mã nguồn, cấu hình và tài
nguyên như lúc lập chỉ mục. Cờ tiền xử lý chính thức xác nhận triển khai đã tích
hợp, không thay cho phép đánh giá chất lượng trên gold segmentation do người gán.

% Phần M1 cho Chương 3 — được \input bên trong chapter3_ir.tex nên KHÔNG có \chapter.
% Chỉ dùng các gói đã khai báo trong main.tex; không thêm gói mới.
\section{Thiết kế tầng ngôn ngữ M1}
\label{sec:m1-thiet-ke}

Tầng ngôn ngữ M1 biến văn bản thô thành dãy token sạch để M2 lập chỉ mục. Tầng này
gồm ba khối độc lập: chuẩn hóa chuỗi, tách từ theo so khớp cực đại trên từ điển, và
lọc từ dừng tầng~1 kèm danh sách bảo vệ. Toàn bộ tầng hoạt động ngoại tuyến, không
gọi dịch vụ ngoài, và mọi tài nguyên mà nó dùng đều được ghim bằng hàm băm.

\subsection{Chuẩn hóa văn bản và bảo toàn ranh giới}

Hàm chuẩn hóa \texttt{normalize(text)} trả về dạng hiển thị: chuẩn hóa Unicode
NFC, hạ chữ thường, loại bỏ thẻ đánh dấu, đường dẫn và địa chỉ thư điện tử, giữ
lại chữ cái tiếng Việt cùng khoảng trắng, và thu gọn mọi khoảng trắng về một dấu
cách. Một ví dụ:

\begin{quote}
\texttt{"Tri thức, nhân loại! 2026"} $\rightarrow$ \texttt{"tri thức nhân loại"}.
\end{quote}

Một quyết định thiết kế quan trọng là \emph{dấu câu là ranh giới, không phải
nhiễu}. Dạng hiển thị gộp khoảng trắng nên trông như một chuỗi liên tục, nhưng bộ
tách từ phải không bao giờ ghép một từ xuyên qua dấu câu. Vì lý do đó, bộ tách từ
tiêu thụ một dạng trung gian giữ lại ranh giới, ở đó đoạn văn bản bị cắt tại mỗi
dấu câu và tại mỗi khoảng trắng không phải dấu cách thường (xuống dòng, tab). Các
khoảng được gỡ bỏ (đường dẫn, thư điện tử, thẻ đánh dấu) cũng được thay bằng một
ranh giới cứng chứ không phải bằng dấu cách; nếu thay bằng dấu cách thì hai mảnh
văn bản hai bên sẽ bị dính thành một từ, ví dụ \texttt{"học<br>sinh"} phải tách
thành hai âm tiết chứ không được ghép thành một từ.

Cũng trong khối này có một chỉnh sửa so với mẫu biểu thức chính quy ban đầu. Mẫu
cũ chỉ giữ chữ cái tiếng Việt nên đã vô tình xóa mọi chữ cái Latin, biến
``Windows'' thành ``indows'' và ``Java'' thành ``ava''. Bốn chữ cái Latin
\texttt{f}, \texttt{j}, \texttt{w}, \texttt{z} được giữ lại vì chúng phân biệt
trong các tên vay mượn phổ biến (Windows, WWW, Java, Zalo). Chúng \emph{không}
được thêm vào từ điển; kết quả trung thực là chúng trở thành token đơn ngoài từ
điển thay vì bị bóp méo thành từ sai.

\subsection{Từ điển và cây Trie}

Tách từ dựa trên từ điển Viet74K \cite{viet74k} gồm 69.858 mục, được ghim hàm băm
SHA256 \texttt{175e8b1b...d8ab} và đi kèm thông tin giấy phép cùng nguồn gốc trong
kho. Từ điển được nạp vào một cây Trie hỗ trợ truy vấn theo tiền tố, cho phép bộ
tách từ kiểm tra đoạn ứng viên có trong từ điển hay không mà không phải so khớp
toàn bộ danh sách. Cửa sổ trượt mặc định lấy bằng độ dài mục dài nhất trong từ
điển, nên không bỏ sót từ nào; người dùng vẫn có thể chỉ định một giới hạn cửa sổ
nhỏ hơn khi cần quan sát thuật toán.

\subsection{Thuật toán FMM, BMM và BiMM}

Ba phương pháp dùng chung một cửa sổ trượt nhưng khác hướng duyệt. FMM duyệt tiến,
BMM duyệt lùi. BiMM chạy cả hai rồi chọn theo bộ quy tắc phân xử sau, xét theo
thứ tự:

\begin{enumerate}
  \item ít token hơn;
  \item ít token một âm tiết hơn;
  \item ít token ngoài từ điển hơn;
  \item nếu vẫn hòa, chọn BMM.
\end{enumerate}

Tiêu chí thứ ba là lưới an toàn tường minh: với cài đặt Trie thuần, mọi token
ngoài từ điển đều được phát ra dưới dạng một âm tiết, nên tiêu chí này được suy ra
từ tiêu chí thứ hai; nhóm vẫn giữ nó cùng khóa so sánh đầy đủ để người đọc kiểm
chứng lại quyết định thay vì tin vào phần tóm tắt.

Ví dụ nhập nhằng kinh điển minh họa vì sao cần phân xử. Với câu ``học sinh học
sinh học'' và một từ điển fixture có cả ``học sinh'' lẫn ``sinh học'':

\begin{quote}
FMM: \texttt{học\_sinh, học\_sinh, học}\\
BMM: \texttt{học, sinh\_học, sinh\_học}\\
BiMM: \texttt{học, sinh\_học, sinh\_học}
\end{quote}

Hai phương án có cùng số token và cùng số token một âm tiết, nên hòa và BiMM chọn
BMM. Ví dụ này dùng một từ điển fixture minh bạch, không phụ thuộc từ điển đóng
gói, để phép so sánh là một khẳng định về \emph{thuật toán} chứ không về một từ
điển cụ thể.

Mỗi lần chạy đều sinh một \emph{vết} từng bước. Mỗi dòng vết ghi số thứ tự bước,
hướng duyệt, chỉ số âm tiết bắt đầu và kết thúc (kết thúc mở), đoạn ứng viên, độ
dài cửa sổ, cờ có khớp từ điển hay không, và hành động -- một trong
\texttt{MATCHED}, \texttt{SHRINK} hoặc \texttt{SINGLE}. Vết dùng đúng từ điển và
cửa sổ của lần chạy, nên người đọc có thể tái lập từng quyết định thay vì chỉ thấy
kết quả cuối cùng. Với đầu vào rỗng hoặc chỉ gồm dấu câu, kết quả là dãy token rỗng
và vết rỗng.

\subsection{Lọc từ dừng tầng 1 và danh sách bảo vệ}

Khối thứ ba loại các hư từ ngữ pháp -- đại từ, giới từ, trợ từ, phó từ chỉ mức độ,
liên từ -- vốn đóng góp ít cho phân biệt truy xuất. Hai tài nguyên có nguồn và quy
tắc biên soạn điều khiển bước này: danh sách hư từ ngữ pháp và danh sách bảo vệ
(guard). Tập từ dừng hiệu lực là hiệu của hai danh sách: \emph{guard luôn thắng}.

Quy tắc bảo vệ có phạm vi rộng hơn một so khớp đơn giản: một token được bảo vệ nếu
chính nó là guard, hoặc nếu một thành phần của nó (tách theo dấu gạch nối âm tiết)
cũng là guard. Nhờ đó việc bảo vệ ``không'' cũng giữ nguyên ``không\_thể'', bảo vệ
``vì'' giữ nguyên ``vì\_dầu'', bảo vệ ``chẳng'' giữ nguyên ``chẳng\_hạn''. Điểm
tinh tế là chỉ những thành phần \emph{tự chúng là guard} mới có thể bảo vệ: nếu mọi
âm tiết của mọi guard đều có thể bảo vệ thì ``thể'' sẽ được bảo vệ và ``thể\_thao''
bị bảo vệ nhầm. Chính vì vậy các từ ghép cần bảo vệ vì lý do khác được liệt kê
tường minh trong tệp guard.

Ví dụ minh họa:

\begin{quote}
\texttt{filter\_tier1(["bóng", "đá", "và", "của", "bóng\_đá"])}
$\rightarrow$ \texttt{["bóng", "đá", "bóng\_đá"]}
\end{quote}

Ở đây ``và'' và ``của'' bị loại, còn ``bóng\_đá'' là một từ ghép nội dung được giữ
nguyên. Cùng cơ chế bảo vệ này được tái dùng cho tầng~2 khi TDV chọn stoplist, để
các từ phủ định không bị xóa ở cả hai tầng.

\subsection{Tài nguyên và xuất xứ}

Ba tài nguyên ngôn ngữ học được đóng gói kèm số mục thực tế và hàm băm, để kết quả
tái lập được và để bất kỳ ai cũng kiểm tra lại được dữ liệu đầu vào.

\begin{table}[ht]
\centering
\begin{tabular}{|p{0.30\textwidth}|p{0.38\textwidth}|p{0.18\textwidth}|}
\hline
Tài nguyên & Vai trò & Số mục \\
\hline
\texttt{dict.txt} & Từ điển tách từ (Viet74K) & 69.858 \\
\texttt{grammar\_\allowbreak stopwords.txt} & Hư từ ngữ pháp tầng 1 & 181 \\
\texttt{guard\_\allowbreak whitelist.txt} & Danh sách bảo vệ (guard) & 110 \\
\hline
\end{tabular}
\caption{Ba tài nguyên của tầng ngôn ngữ M1 (đặt trong thư mục \texttt{resources/})
và số mục tại thời điểm viết. Số mục được đọc trực tiếp từ tệp; danh sách hư từ sau
khi trừ guard còn 124 từ dừng hiệu lực.}
\label{tab:m1-tai-nguyen}
\end{table}

Từ điển là danh sách mã nguồn mở có ghi giấy phép \cite{viet74k}. Danh sách từ
dừng tĩnh dùng cho cấu hình C2 là một bản chụp có nguồn, phiên bản và thông báo
giấy phép MIT \cite{static_stopwords}. Hai danh sách tầng~1 do nhóm biên soạn kèm
quy tắc và số mục thực tế; các từ dừng được gỡ lại chỉ là các từ phủ định và logic
có lý do ngữ nghĩa rõ ràng, không bù danh sách bằng thực từ để đạt một con số.

\subsection{Kiểm chứng tầng ngôn ngữ}

Ba phương pháp tách từ được kiểm chứng bằng từ điển fixture minh bạch và bằng một
đường tham chiếu trực tiếp trên từ điển/câu ngẫu nhiên. Vết thuật toán cho biết ứng
viên và quyết định tại từng bước; dãy token do \texttt{tokenize} trả về được kiểm
riêng bằng fixture, vì với BMM các dòng ứng viên được đánh dấu có thể chồng lấn nên
không dùng chúng để tái dựng kết quả. Bộ lọc tầng~1 được kiểm chứng
bằng các ca guard: phủ định và từ ghép chứa guard phải được giữ nguyên, còn hư từ
không nằm trong guard phải bị loại. Các ca kiểm thử này chứng minh \emph{luật hoạt
động đúng}, không đo \emph{chất lượng} tách từ.

Chất lượng tách từ chỉ đo được khi có nhãn người. Nhóm đã dựng bộ mẫu 150 câu
(15 câu mỗi lĩnh vực) rút từ kho ngữ liệu đã khóa, kèm công cụ gán nhãn ngoại
tuyến và bộ đo khớp span từ. Hiện 150/150 câu có cờ xác nhận và phép đo word-span
đã chạy, nhưng ba câu SG010, SG053, SG089 còn cần rà lại về ranh giới từ; do đó
F1 ở Chương~4 là số đo tạm thời trên gold hiện có, không phải kết luận chất lượng
cuối. Gợi ý của máy chỉ phục vụ gán nhãn nhanh và không tự động trở thành gold.
Quy trình đo, các ràng buộc theo hàm băm và giới hạn của phép đo trình bày trong
Chương~4.


\section{Độ đo TF-IDF và chuẩn hóa}

Với tài liệu $d$ và từ $t$:

\begin{align}
TF(t,d) &= 1+\ln f_{t,d} \quad (f_{t,d}>0), \\
IDF(t) &= \ln\frac{N+1}{df(t)+1}+1, \\
w_{t,d} &= TF(t,d)\times IDF(t), \\
\vec{x}_d &= \frac{w_d}{\lVert w_d\rVert_2}.
\end{align}

Ba lựa chọn cần nêu rõ vì chúng ảnh hưởng trực tiếp tới kết quả đo. Thứ nhất,
dùng sublinear TF nghĩa là một từ xuất hiện ba lần không được coi là gấp ba lần
một từ xuất hiện một lần. Thứ hai, IDF làm mượt $\ln\frac{N+1}{df+1}+1$ không bao
giờ bằng không, nên không mất thông tin ở từ xuất hiện trong mọi tài liệu. Thứ
ba, chuẩn L2 biến tích số vào thành cosin, cho phép so sánh vector có độ dài
khác nhau.

\section{Độ đo TDV và lọc stopword tầng 2}

\subsection{Vì sao cần TDV}

Bộ stopword tĩnh C2 là danh sách tiếng Việt có nguồn và phiên bản cố định. Với
tin tức tiếng Việt theo mười chủ đề, có những từ rất phổ biến nhưng lại chỉ xuất
hiện trong một vài chủ đề; loại chúng đi có thể làm mất thông tin phân biệt. TDV
(Salton) đo trực tiếp câu hỏi đó: bỏ một từ khỏi không gian thì các tài liệu có
dễ phân biệt nhau hơn hay không?

\subsection{Định nghĩa mật độ}

Với $N$ tài liệu đã chuẩn hóa L2, mật độ không gian tài liệu là mật độ cosin
pairwise, loại trừ self-pair:

\begin{equation}
D(X)=\frac{1}{N(N-1)}\sum_{i\neq j}\cos(x_i,x_j).
\end{equation}

Mẫu số $N(N-1)$ là số cặp có thứ tự, được giữ cố định bất kể có bao nhiêu tài
liệu thực sự khác nhau. Với $N<2$ ta đặt $D=0$. Vectơ tất cả bằng không thì
$\cos=0$ với mọi tài liệu khác, nên nó đóng góp bằng không cho mọi cặp chứ không
bị loại khỏi mẫu số.

\subsection{Định nghĩa TDV}

\begin{equation}
TDV(t)=D_{-t}-D_{base},
\end{equation}

trong đó $D_{-t}$ là mật độ sau khi xoá hẳn chiều $t$. Điểm dễ sai nhất nằm
đúng ở chỗ này. Xoá một chiều làm thay đổi hai thứ cùng lúc: độ dài vector của
từng tài liệu, và do đó toàn bộ tổng vector. Vì vậy sau khi xoá chiều, mỗi tài
liệu phải được chuẩn hóa L2 lại, và tổng vector phải được tính lại trên tập
vector đã chuẩn hóa mới.

\subsection{Trọng số đo TDV và tính một lần duy nhất}

TDV được đo trên TF-IDF sơ bộ, tức là trước khi áp bất kỳ bộ lọc TDV nào, và
\textbf{giữ nguyên IDF sơ bộ} khi loại một chiều. Nếu tính lại IDF sau mỗi lần
loại thì hai hiệu ứng trộn vào nhau và con số thu được không còn là TDV của riêng
từ đang xét.

Điểm thứ hai cần nhấn mạnh là TDV được tính \textbf{một lần duy nhất trên toàn
bộ corpus}: mọi từ cùng chia sẻ một ma trận gốc. Không lấy điểm của từ đã bị loại
ở bước trước rồi mới xoá thêm từ sau, vì khi đó điểm đo phụ thuộc thứ tự xoá và
không còn là thuộc tính của từ.

\subsection{Cách cài đặt và kiểm chứng}

Gọi $v_i$ là thành phần chiều $t$ của tài liệu $i$, $S=X X^{\mathsf T}$. Sau khi
xoá chiều $t$, vector của tài liệu $i$ là $x_i - v_i e_t$ với độ dài
$n_i=\sqrt{\lVert x_i\rVert^2-v_i^2}$. Đặt $u_i=1/n_i$, và $u_i=0$ khi vectơ sau
khi xoá bằng không. Khi đó tổng cosin ngoài đường chéo sau khi xoá là

\begin{equation}
\sum_{i\neq j} u_iu_j\,(S_{ij}-v_iv_j)
= \lVert X^{\mathsf T}u\rVert^2 - \sum_i u_i^2 S_{ii}
- \Bigl(\langle u,v\rangle^2-\sum_i u_i^2 v_i^2\Bigr).
\end{equation}

Nhờ đẳng thức này, mỗi từ chỉ cần duyệt posting list của nó thay vì dựng lại ma
trận cosin đầy đủ $N\times N$.

Cách làm sai mà đề cương cảnh báo, tức là xoá chiều mà không tái chuẩn hóa và
không cập nhật tổng vector, được cài sẵn trong bộ kiểm thử dưới dạng một hàm tham
chiếu thứ hai. Bộ test khẳng định kết quả của engine khác kết quả của cách làm
sai, đồng thời khớp với một hàm tham chiếu thứ ba hoàn toàn trực tiếp: xoá cột,
tái chuẩn hóa từng dòng, nhân ma trận cosin đầy đủ.

\subsection{Ngưỡng lọc}

\begin{equation}
\text{giữ từ } t \iff TDV(t)\geq -\varepsilon,\qquad \varepsilon=10^{-12}.
\end{equation}

$\varepsilon$ là dung sai số học chứ không phải siêu tham số thuật toán truy xuất;
giá trị này được ghi vào metadata của mọi lần chạy. Không ép đường cong hình chữ
U, và không suy ra từ dấu của TDV rằng một cấu hình lọc nào sẽ cho kết quả tốt
hơn: việc loại một từ có cải thiện truy xuất hay không chỉ trả lời được bằng phép
đo với qrels đã gán nhãn.

\section{Bốn cấu hình ablation}

\begin{center}
\begin{tabular}{|l|l|l|}
\hline
Mã & Cấu hình & Ghi chú \\
\hline
C1 & Không lọc stopword & baseline \\
C2 & Stoplist tĩnh, độc lập với corpus & nguồn và số mục thực tế được ghi kèm \\
C3 & Chỉ stopword ngữ pháp tầng 1 của M1 & \\
C4 & Tầng 1 cộng TDV tầng 2 & giữ các từ phủ định bằng guard ngữ nghĩa \\
\hline
\end{tabular}
\end{center}

Mỗi cấu hình nằm trong một thư mục riêng kèm metadata: SHA256 corpus, tokenizer,
cấu hình, stoplist và số phần tử stoplist. Bốn chỉ mục dùng chung một corpus, một
bộ truy vấn và một bộ nhãn, nên so sánh giữa chúng là so sánh đúng yếu tố lọc
stopword.

\subsection{Về lớp tầng 1 tạm thời}

Hợp đồng dự án quy định: vì đồ án chưa có Maximum Matching chính thức của M1, nếu
C3/C4 dùng bộ lọc tầng 1 tạm thì kết quả phải được ghi là provisional, không
coi là Maximum Matching chính thức. Bộ đánh giá vì vậy chặn tuyệt đối mọi kết quả
chính thức khi metadata còn trỏ tới bộ lọc tạm hoặc tokenizer phát triển, kể cả
khi metadata ghi nhầm trạng thái là official.

\subsection{Giữ từ phủ định}

Bộ TDV theo bản chất thiên về việc loại từ phổ biến, mà các từ phủ định như
``không'', ``chưa'', ``chẳng'', ``vì'' lại mang thông tin quan trọng cho truy vấn
dạng phủ định. Vì vậy sau khi TDV chọn stoplist, một danh sách guard ngữ nghĩa
được gỡ lại khỏi stoplist. Với corpus chính, TDV chọn 498 từ và sau bước guard còn
494 từ, với bốn từ được giữ lại là ``chưa'', ``chẳng'', ``không'' và ``vì''.

\subsection{Cơ chế fallback của C4 và giới hạn của nó}

Khi một truy vấn mà mọi token đều đã bị lọc khỏi từ vựng của C4, vectơ truy vấn
bằng không và không thể có kết quả. Cơ chế dự phòng là tra cứu bằng một chỉ mục
C3 riêng, có cùng corpus và cùng tiền xử lý. Cần nói rõ hai điều. Một, giữ lại một
từ đã bị xoá khỏi chỉ mục không phục hồi được khả năng truy xuất bằng chỉ mục đó;
fallback chỉ cứu được truy vấn mà không còn token nào. Hai, hệ thống khi đó phải
giữ cả hai chỉ mục, nên dung lượng triển khai tăng lên chứ không giảm; số liệu
mục 6.1 của Chương 4 đo riêng cả chỉ mục C4 thuần lẫn gói C4 cộng C3.

\section{Cấu trúc chỉ mục và đường đi của một truy vấn}

\begin{enumerate}
\item Chuẩn hóa văn bản và tách từ theo đối tượng tiền xử lý của M1.
\item Lọc stopword theo đúng cấu hình đã dùng khi lập chỉ mục.
\item Vector hoá truy vấn bằng bảng IDF ngoại tuyến của corpus.
\item Lấy ứng viên bằng hợp các posting list trong chỉ mục đảo.
\item Tính cosin trên tập ứng viên, chấm điểm chỉ những kết quả có điểm dương.
\item Sắp xếp giảm dần điểm; hòa điểm sắp theo \texttt{doc\_id} tăng dần để kết
quả tái lập được.
\end{enumerate}

Chỉ mục được lưu dạng thưa: posting list và vector trọng số. Hàm giải thích
(\texttt{explain}) trả về từng đóng góp vào tích vô hướng, và tổng các đóng góp
bằng đúng điểm cosin đã trả về, để người đọc có thể kiểm tra lại cách xếp hạng.

\section{Giao diện đánh giá và cổng chặn kết quả chính thức}

Bộ đánh giá dùng các định nghĩa chỉ số sau, với $K$ mặc định là 10 và có thể cấu
hình:

\begin{align}
P@K &= \frac{\text{số trúng trong Top-}K}{K}, \\
Recall@K &= \frac{\text{số trúng trong Top-}K}{\text{tổng số relevant đã confirmed}}, \\
F1@K &= \frac{2 \cdot P@K \cdot Recall@K}{P@K + Recall@K}, \\
AP &= \frac{1}{R}\sum_{i=1}^{L} P@i \cdot rel_i, \\
NDCG@K &= \frac{DCG}{IDCG},\quad DCG=\sum_{i=1}^{K}\frac{2^{rel_i}-1}{\log_2(i+1)} .
\end{align}

Ba quyết định thiết kế đáng ghi lại. Mẫu số của $P@K$ là $K$ cố định, kể cả khi
hệ thống trả về ít hơn $K$ kết quả; thiếu kết quả là sai chứ không phải là không
tính. $AP$ tính trên toàn bộ ranking có sẵn, không cắt theo $K$, nên bộ đánh giá
từ chối chạy nếu độ sâu lấy ranking nhỏ hơn số tài liệu corpus. $IDCG$ lấy từ
qrels chứ không lấy từ lần chạy, nên điểm không thể bị đẩy lên bằng cách chỉ
trả về một phần kết quả.

Query không có tài liệu relevant nào vẫn được giữ trong bình quân với giá trị 0 và
được đếm riêng, để người đọc biết chính xác MAP bị kéo xuống bao nhiêu do những
query đó. Tương tự, một truy vấn không có kết quả trả về vẫn sinh ra dòng chỉ số
bằng 0 chứ không bị loại khỏi bình quân.

\subsection{Điều kiện để coi là đánh giá chính thức}

Bộ đánh gia từ chối ghi kết quả chính thức khi có bất kỳ điều kiện nào sau đây:

\begin{itemize}
\item dưới 20 truy vấn;
\item SHA256 corpus lệch giữa manifest, qrels, tập truy vấn hoặc chỉ mục, hoặc
thiếu hẳn hash ở bất kỳ nguồn nào;
\item còn phán quyết chưa confirmed, hoặc nhãn confirmed mà điểm relevance còn
null, hoặc nhãn confirmed không ghi người gán;
\item trùng cặp (truy vấn, tài liệu), điểm relevance ngoài tập $\{0,1,2\}$, hoặc
truy vấn trong qrels không tồn tại trong tập truy vấn;
\item metadata tiền xử lý chưa nhất quán ở trạng thái official, hoặc đang dùng
tokenizer phát triển hay bộ lọc tầng 1 tạm.
\end{itemize}

SHA256 corpus cần đánh giá không bao giờ được lấy từ qrels: nếu không xác định
được từ manifest đã khóa hoặc từ chính file corpus thì kết quả chính thức bị chặn.

\subsection{Khi chưa đủ nhãn thì không xuất chỉ số}

Nếu bất kỳ truy vấn nào chưa có nhãn confirmed, bộ đánh giá ghi
\texttt{metrics\_status = unmeasured} và không ghi bảng chỉ số. Lý do rất cụ thể:
với qrels rỗng, mọi chỉ số đều bằng không một cách máy móc, và một bảng toàn số
không dễ bị đọc nhầm là hệ thống không tìm được gì.

\section{Tách kết quả kiểm thử khỏi kết quả chính thức}

Các bài kiểm thử dùng dữ liệu tạm trong thư mục làm việc của pytest. Để ngăn
tình huống số liệu kiểm thử lọt vào báo cáo, bộ đánh giá chặn ghi kết quả không
chính thức vào các thư mục kết quả chính thức; kết quả chạy thử bắt buộc ghi ra
thư mục khác, ví dụ \texttt{results/development\_eval}.
 trong main.tex.
% Chỉ dùng các gói đã khai báo trong main.tex: report, vietnamese, amsmath,
% amssymb, graphicx, hyperref, geometry. Không thêm gói mới.
\chapter{THIẾT KẾ HỆ THỐNG TÌM KIẾM VĂN BẢN}

\section{Phạm vi trách nhiệm của M2 trong kiến trúc}

Kiến trúc tổng thể chia thành ba lớp trách nhiệm. M1 sở hữu tiền xử lý ngôn ngữ
tự nhiên, M2 sở hữu biến các từ đã tách thành dữ liệu lập chỉ mục và truy xuất,
M3 sở hữu giao diện người dùng. Chương này trình bày phần của M2.

Các lớp và ranh giới tích hợp được trình bày trong Bảng~\ref{tab:kien-truc-m2}.
\begin{table}[ht]
\centering
\begin{tabular}{|p{0.08\textwidth}|p{0.43\textwidth}|p{0.35\textwidth}|}
\hline
Lớp & Thành phần & File chính \\
\hline
M1 & Maximum Matching, stopword tầng 1 & \texttt{engine/tokenizer/} \\
M2 & Adapter M1 có fingerprint & \texttt{engine/preprocessing.py} \\
M2 & TDV, TF-IDF, chỉ mục đảo, tìm kiếm & \texttt{engine/indexer/} \\
M2 & Đánh giá, đo độ trễ, biểu đồ & \texttt{evaluator/}, \texttt{benchmarks/} \\
M3 & Giao diện Streamlit & \texttt{app/} \\
\hline
\end{tabular}
\caption{Các lớp chức năng và ranh giới tích hợp của M2.}
\label{tab:kien-truc-m2}
\end{table}

Phiên bản tích hợp nhận BiMM và bộ lọc tầng 1 từ M1 qua hai callable độc lập.
Baseline âm tiết và fixture tầng 1 chỉ còn dùng cho kiểm thử/phát triển,
được đánh dấu \texttt{development\_only}. Khi nạp chỉ mục,
đối tượng tiền xử lý phải có cùng fingerprint của mã nguồn, cấu hình và tài
nguyên như lúc lập chỉ mục. Cờ tiền xử lý chính thức xác nhận triển khai đã tích
hợp, không thay cho phép đánh giá chất lượng trên gold segmentation do người gán.

% Phần M1 cho Chương 3 — được \input bên trong chapter3_ir.tex nên KHÔNG có \chapter.
% Chỉ dùng các gói đã khai báo trong main.tex; không thêm gói mới.
\section{Thiết kế tầng ngôn ngữ M1}
\label{sec:m1-thiet-ke}

Tầng ngôn ngữ M1 biến văn bản thô thành dãy token sạch để M2 lập chỉ mục. Tầng này
gồm ba khối độc lập: chuẩn hóa chuỗi, tách từ theo so khớp cực đại trên từ điển, và
lọc từ dừng tầng~1 kèm danh sách bảo vệ. Toàn bộ tầng hoạt động ngoại tuyến, không
gọi dịch vụ ngoài, và mọi tài nguyên mà nó dùng đều được ghim bằng hàm băm.

\subsection{Chuẩn hóa văn bản và bảo toàn ranh giới}

Hàm chuẩn hóa \texttt{normalize(text)} trả về dạng hiển thị: chuẩn hóa Unicode
NFC, hạ chữ thường, loại bỏ thẻ đánh dấu, đường dẫn và địa chỉ thư điện tử, giữ
lại chữ cái tiếng Việt cùng khoảng trắng, và thu gọn mọi khoảng trắng về một dấu
cách. Một ví dụ:

\begin{quote}
\texttt{"Tri thức, nhân loại! 2026"} $\rightarrow$ \texttt{"tri thức nhân loại"}.
\end{quote}

Một quyết định thiết kế quan trọng là \emph{dấu câu là ranh giới, không phải
nhiễu}. Dạng hiển thị gộp khoảng trắng nên trông như một chuỗi liên tục, nhưng bộ
tách từ phải không bao giờ ghép một từ xuyên qua dấu câu. Vì lý do đó, bộ tách từ
tiêu thụ một dạng trung gian giữ lại ranh giới, ở đó đoạn văn bản bị cắt tại mỗi
dấu câu và tại mỗi khoảng trắng không phải dấu cách thường (xuống dòng, tab). Các
khoảng được gỡ bỏ (đường dẫn, thư điện tử, thẻ đánh dấu) cũng được thay bằng một
ranh giới cứng chứ không phải bằng dấu cách; nếu thay bằng dấu cách thì hai mảnh
văn bản hai bên sẽ bị dính thành một từ, ví dụ \texttt{"học<br>sinh"} phải tách
thành hai âm tiết chứ không được ghép thành một từ.

Cũng trong khối này có một chỉnh sửa so với mẫu biểu thức chính quy ban đầu. Mẫu
cũ chỉ giữ chữ cái tiếng Việt nên đã vô tình xóa mọi chữ cái Latin, biến
``Windows'' thành ``indows'' và ``Java'' thành ``ava''. Bốn chữ cái Latin
\texttt{f}, \texttt{j}, \texttt{w}, \texttt{z} được giữ lại vì chúng phân biệt
trong các tên vay mượn phổ biến (Windows, WWW, Java, Zalo). Chúng \emph{không}
được thêm vào từ điển; kết quả trung thực là chúng trở thành token đơn ngoài từ
điển thay vì bị bóp méo thành từ sai.

\subsection{Từ điển và cây Trie}

Tách từ dựa trên từ điển Viet74K \cite{viet74k} gồm 69.858 mục, được ghim hàm băm
SHA256 \texttt{175e8b1b...d8ab} và đi kèm thông tin giấy phép cùng nguồn gốc trong
kho. Từ điển được nạp vào một cây Trie hỗ trợ truy vấn theo tiền tố, cho phép bộ
tách từ kiểm tra đoạn ứng viên có trong từ điển hay không mà không phải so khớp
toàn bộ danh sách. Cửa sổ trượt mặc định lấy bằng độ dài mục dài nhất trong từ
điển, nên không bỏ sót từ nào; người dùng vẫn có thể chỉ định một giới hạn cửa sổ
nhỏ hơn khi cần quan sát thuật toán.

\subsection{Thuật toán FMM, BMM và BiMM}

Ba phương pháp dùng chung một cửa sổ trượt nhưng khác hướng duyệt. FMM duyệt tiến,
BMM duyệt lùi. BiMM chạy cả hai rồi chọn theo bộ quy tắc phân xử sau, xét theo
thứ tự:

\begin{enumerate}
  \item ít token hơn;
  \item ít token một âm tiết hơn;
  \item ít token ngoài từ điển hơn;
  \item nếu vẫn hòa, chọn BMM.
\end{enumerate}

Tiêu chí thứ ba là lưới an toàn tường minh: với cài đặt Trie thuần, mọi token
ngoài từ điển đều được phát ra dưới dạng một âm tiết, nên tiêu chí này được suy ra
từ tiêu chí thứ hai; nhóm vẫn giữ nó cùng khóa so sánh đầy đủ để người đọc kiểm
chứng lại quyết định thay vì tin vào phần tóm tắt.

Ví dụ nhập nhằng kinh điển minh họa vì sao cần phân xử. Với câu ``học sinh học
sinh học'' và một từ điển fixture có cả ``học sinh'' lẫn ``sinh học'':

\begin{quote}
FMM: \texttt{học\_sinh, học\_sinh, học}\\
BMM: \texttt{học, sinh\_học, sinh\_học}\\
BiMM: \texttt{học, sinh\_học, sinh\_học}
\end{quote}

Hai phương án có cùng số token và cùng số token một âm tiết, nên hòa và BiMM chọn
BMM. Ví dụ này dùng một từ điển fixture minh bạch, không phụ thuộc từ điển đóng
gói, để phép so sánh là một khẳng định về \emph{thuật toán} chứ không về một từ
điển cụ thể.

Mỗi lần chạy đều sinh một \emph{vết} từng bước. Mỗi dòng vết ghi số thứ tự bước,
hướng duyệt, chỉ số âm tiết bắt đầu và kết thúc (kết thúc mở), đoạn ứng viên, độ
dài cửa sổ, cờ có khớp từ điển hay không, và hành động -- một trong
\texttt{MATCHED}, \texttt{SHRINK} hoặc \texttt{SINGLE}. Vết dùng đúng từ điển và
cửa sổ của lần chạy, nên người đọc có thể tái lập từng quyết định thay vì chỉ thấy
kết quả cuối cùng. Với đầu vào rỗng hoặc chỉ gồm dấu câu, kết quả là dãy token rỗng
và vết rỗng.

\subsection{Lọc từ dừng tầng 1 và danh sách bảo vệ}

Khối thứ ba loại các hư từ ngữ pháp -- đại từ, giới từ, trợ từ, phó từ chỉ mức độ,
liên từ -- vốn đóng góp ít cho phân biệt truy xuất. Hai tài nguyên có nguồn và quy
tắc biên soạn điều khiển bước này: danh sách hư từ ngữ pháp và danh sách bảo vệ
(guard). Tập từ dừng hiệu lực là hiệu của hai danh sách: \emph{guard luôn thắng}.

Quy tắc bảo vệ có phạm vi rộng hơn một so khớp đơn giản: một token được bảo vệ nếu
chính nó là guard, hoặc nếu một thành phần của nó (tách theo dấu gạch nối âm tiết)
cũng là guard. Nhờ đó việc bảo vệ ``không'' cũng giữ nguyên ``không\_thể'', bảo vệ
``vì'' giữ nguyên ``vì\_dầu'', bảo vệ ``chẳng'' giữ nguyên ``chẳng\_hạn''. Điểm
tinh tế là chỉ những thành phần \emph{tự chúng là guard} mới có thể bảo vệ: nếu mọi
âm tiết của mọi guard đều có thể bảo vệ thì ``thể'' sẽ được bảo vệ và ``thể\_thao''
bị bảo vệ nhầm. Chính vì vậy các từ ghép cần bảo vệ vì lý do khác được liệt kê
tường minh trong tệp guard.

Ví dụ minh họa:

\begin{quote}
\texttt{filter\_tier1(["bóng", "đá", "và", "của", "bóng\_đá"])}
$\rightarrow$ \texttt{["bóng", "đá", "bóng\_đá"]}
\end{quote}

Ở đây ``và'' và ``của'' bị loại, còn ``bóng\_đá'' là một từ ghép nội dung được giữ
nguyên. Cùng cơ chế bảo vệ này được tái dùng cho tầng~2 khi TDV chọn stoplist, để
các từ phủ định không bị xóa ở cả hai tầng.

\subsection{Tài nguyên và xuất xứ}

Ba tài nguyên ngôn ngữ học được đóng gói kèm số mục thực tế và hàm băm, để kết quả
tái lập được và để bất kỳ ai cũng kiểm tra lại được dữ liệu đầu vào.

\begin{table}[ht]
\centering
\begin{tabular}{|p{0.30\textwidth}|p{0.38\textwidth}|p{0.18\textwidth}|}
\hline
Tài nguyên & Vai trò & Số mục \\
\hline
\texttt{dict.txt} & Từ điển tách từ (Viet74K) & 69.858 \\
\texttt{grammar\_\allowbreak stopwords.txt} & Hư từ ngữ pháp tầng 1 & 181 \\
\texttt{guard\_\allowbreak whitelist.txt} & Danh sách bảo vệ (guard) & 110 \\
\hline
\end{tabular}
\caption{Ba tài nguyên của tầng ngôn ngữ M1 (đặt trong thư mục \texttt{resources/})
và số mục tại thời điểm viết. Số mục được đọc trực tiếp từ tệp; danh sách hư từ sau
khi trừ guard còn 124 từ dừng hiệu lực.}
\label{tab:m1-tai-nguyen}
\end{table}

Từ điển là danh sách mã nguồn mở có ghi giấy phép \cite{viet74k}. Danh sách từ
dừng tĩnh dùng cho cấu hình C2 là một bản chụp có nguồn, phiên bản và thông báo
giấy phép MIT \cite{static_stopwords}. Hai danh sách tầng~1 do nhóm biên soạn kèm
quy tắc và số mục thực tế; các từ dừng được gỡ lại chỉ là các từ phủ định và logic
có lý do ngữ nghĩa rõ ràng, không bù danh sách bằng thực từ để đạt một con số.

\subsection{Kiểm chứng tầng ngôn ngữ}

Ba phương pháp tách từ được kiểm chứng bằng từ điển fixture minh bạch và bằng một
đường tham chiếu trực tiếp trên từ điển/câu ngẫu nhiên. Vết thuật toán cho biết ứng
viên và quyết định tại từng bước; dãy token do \texttt{tokenize} trả về được kiểm
riêng bằng fixture, vì với BMM các dòng ứng viên được đánh dấu có thể chồng lấn nên
không dùng chúng để tái dựng kết quả. Bộ lọc tầng~1 được kiểm chứng
bằng các ca guard: phủ định và từ ghép chứa guard phải được giữ nguyên, còn hư từ
không nằm trong guard phải bị loại. Các ca kiểm thử này chứng minh \emph{luật hoạt
động đúng}, không đo \emph{chất lượng} tách từ.

Chất lượng tách từ chỉ đo được khi có nhãn người. Nhóm đã dựng bộ mẫu 150 câu
(15 câu mỗi lĩnh vực) rút từ kho ngữ liệu đã khóa, kèm công cụ gán nhãn ngoại
tuyến và bộ đo khớp span từ. Hiện 150/150 câu có cờ xác nhận và phép đo word-span
đã chạy, nhưng ba câu SG010, SG053, SG089 còn cần rà lại về ranh giới từ; do đó
F1 ở Chương~4 là số đo tạm thời trên gold hiện có, không phải kết luận chất lượng
cuối. Gợi ý của máy chỉ phục vụ gán nhãn nhanh và không tự động trở thành gold.
Quy trình đo, các ràng buộc theo hàm băm và giới hạn của phép đo trình bày trong
Chương~4.


\section{Độ đo TF-IDF và chuẩn hóa}

Với tài liệu $d$ và từ $t$:

\begin{align}
TF(t,d) &= 1+\ln f_{t,d} \quad (f_{t,d}>0), \\
IDF(t) &= \ln\frac{N+1}{df(t)+1}+1, \\
w_{t,d} &= TF(t,d)\times IDF(t), \\
\vec{x}_d &= \frac{w_d}{\lVert w_d\rVert_2}.
\end{align}

Ba lựa chọn cần nêu rõ vì chúng ảnh hưởng trực tiếp tới kết quả đo. Thứ nhất,
dùng sublinear TF nghĩa là một từ xuất hiện ba lần không được coi là gấp ba lần
một từ xuất hiện một lần. Thứ hai, IDF làm mượt $\ln\frac{N+1}{df+1}+1$ không bao
giờ bằng không, nên không mất thông tin ở từ xuất hiện trong mọi tài liệu. Thứ
ba, chuẩn L2 biến tích số vào thành cosin, cho phép so sánh vector có độ dài
khác nhau.

\section{Độ đo TDV và lọc stopword tầng 2}

\subsection{Vì sao cần TDV}

Bộ stopword tĩnh C2 là danh sách tiếng Việt có nguồn và phiên bản cố định. Với
tin tức tiếng Việt theo mười chủ đề, có những từ rất phổ biến nhưng lại chỉ xuất
hiện trong một vài chủ đề; loại chúng đi có thể làm mất thông tin phân biệt. TDV
(Salton) đo trực tiếp câu hỏi đó: bỏ một từ khỏi không gian thì các tài liệu có
dễ phân biệt nhau hơn hay không?

\subsection{Định nghĩa mật độ}

Với $N$ tài liệu đã chuẩn hóa L2, mật độ không gian tài liệu là mật độ cosin
pairwise, loại trừ self-pair:

\begin{equation}
D(X)=\frac{1}{N(N-1)}\sum_{i\neq j}\cos(x_i,x_j).
\end{equation}

Mẫu số $N(N-1)$ là số cặp có thứ tự, được giữ cố định bất kể có bao nhiêu tài
liệu thực sự khác nhau. Với $N<2$ ta đặt $D=0$. Vectơ tất cả bằng không thì
$\cos=0$ với mọi tài liệu khác, nên nó đóng góp bằng không cho mọi cặp chứ không
bị loại khỏi mẫu số.

\subsection{Định nghĩa TDV}

\begin{equation}
TDV(t)=D_{-t}-D_{base},
\end{equation}

trong đó $D_{-t}$ là mật độ sau khi xoá hẳn chiều $t$. Điểm dễ sai nhất nằm
đúng ở chỗ này. Xoá một chiều làm thay đổi hai thứ cùng lúc: độ dài vector của
từng tài liệu, và do đó toàn bộ tổng vector. Vì vậy sau khi xoá chiều, mỗi tài
liệu phải được chuẩn hóa L2 lại, và tổng vector phải được tính lại trên tập
vector đã chuẩn hóa mới.

\subsection{Trọng số đo TDV và tính một lần duy nhất}

TDV được đo trên TF-IDF sơ bộ, tức là trước khi áp bất kỳ bộ lọc TDV nào, và
\textbf{giữ nguyên IDF sơ bộ} khi loại một chiều. Nếu tính lại IDF sau mỗi lần
loại thì hai hiệu ứng trộn vào nhau và con số thu được không còn là TDV của riêng
từ đang xét.

Điểm thứ hai cần nhấn mạnh là TDV được tính \textbf{một lần duy nhất trên toàn
bộ corpus}: mọi từ cùng chia sẻ một ma trận gốc. Không lấy điểm của từ đã bị loại
ở bước trước rồi mới xoá thêm từ sau, vì khi đó điểm đo phụ thuộc thứ tự xoá và
không còn là thuộc tính của từ.

\subsection{Cách cài đặt và kiểm chứng}

Gọi $v_i$ là thành phần chiều $t$ của tài liệu $i$, $S=X X^{\mathsf T}$. Sau khi
xoá chiều $t$, vector của tài liệu $i$ là $x_i - v_i e_t$ với độ dài
$n_i=\sqrt{\lVert x_i\rVert^2-v_i^2}$. Đặt $u_i=1/n_i$, và $u_i=0$ khi vectơ sau
khi xoá bằng không. Khi đó tổng cosin ngoài đường chéo sau khi xoá là

\begin{equation}
\sum_{i\neq j} u_iu_j\,(S_{ij}-v_iv_j)
= \lVert X^{\mathsf T}u\rVert^2 - \sum_i u_i^2 S_{ii}
- \Bigl(\langle u,v\rangle^2-\sum_i u_i^2 v_i^2\Bigr).
\end{equation}

Nhờ đẳng thức này, mỗi từ chỉ cần duyệt posting list của nó thay vì dựng lại ma
trận cosin đầy đủ $N\times N$.

Cách làm sai mà đề cương cảnh báo, tức là xoá chiều mà không tái chuẩn hóa và
không cập nhật tổng vector, được cài sẵn trong bộ kiểm thử dưới dạng một hàm tham
chiếu thứ hai. Bộ test khẳng định kết quả của engine khác kết quả của cách làm
sai, đồng thời khớp với một hàm tham chiếu thứ ba hoàn toàn trực tiếp: xoá cột,
tái chuẩn hóa từng dòng, nhân ma trận cosin đầy đủ.

\subsection{Ngưỡng lọc}

\begin{equation}
\text{giữ từ } t \iff TDV(t)\geq -\varepsilon,\qquad \varepsilon=10^{-12}.
\end{equation}

$\varepsilon$ là dung sai số học chứ không phải siêu tham số thuật toán truy xuất;
giá trị này được ghi vào metadata của mọi lần chạy. Không ép đường cong hình chữ
U, và không suy ra từ dấu của TDV rằng một cấu hình lọc nào sẽ cho kết quả tốt
hơn: việc loại một từ có cải thiện truy xuất hay không chỉ trả lời được bằng phép
đo với qrels đã gán nhãn.

\section{Bốn cấu hình ablation}

\begin{center}
\begin{tabular}{|l|l|l|}
\hline
Mã & Cấu hình & Ghi chú \\
\hline
C1 & Không lọc stopword & baseline \\
C2 & Stoplist tĩnh, độc lập với corpus & nguồn và số mục thực tế được ghi kèm \\
C3 & Chỉ stopword ngữ pháp tầng 1 của M1 & \\
C4 & Tầng 1 cộng TDV tầng 2 & giữ các từ phủ định bằng guard ngữ nghĩa \\
\hline
\end{tabular}
\end{center}

Mỗi cấu hình nằm trong một thư mục riêng kèm metadata: SHA256 corpus, tokenizer,
cấu hình, stoplist và số phần tử stoplist. Bốn chỉ mục dùng chung một corpus, một
bộ truy vấn và một bộ nhãn, nên so sánh giữa chúng là so sánh đúng yếu tố lọc
stopword.

\subsection{Về lớp tầng 1 tạm thời}

Hợp đồng dự án quy định: vì đồ án chưa có Maximum Matching chính thức của M1, nếu
C3/C4 dùng bộ lọc tầng 1 tạm thì kết quả phải được ghi là provisional, không
coi là Maximum Matching chính thức. Bộ đánh giá vì vậy chặn tuyệt đối mọi kết quả
chính thức khi metadata còn trỏ tới bộ lọc tạm hoặc tokenizer phát triển, kể cả
khi metadata ghi nhầm trạng thái là official.

\subsection{Giữ từ phủ định}

Bộ TDV theo bản chất thiên về việc loại từ phổ biến, mà các từ phủ định như
``không'', ``chưa'', ``chẳng'', ``vì'' lại mang thông tin quan trọng cho truy vấn
dạng phủ định. Vì vậy sau khi TDV chọn stoplist, một danh sách guard ngữ nghĩa
được gỡ lại khỏi stoplist. Với corpus chính, TDV chọn 498 từ và sau bước guard còn
494 từ, với bốn từ được giữ lại là ``chưa'', ``chẳng'', ``không'' và ``vì''.

\subsection{Cơ chế fallback của C4 và giới hạn của nó}

Khi một truy vấn mà mọi token đều đã bị lọc khỏi từ vựng của C4, vectơ truy vấn
bằng không và không thể có kết quả. Cơ chế dự phòng là tra cứu bằng một chỉ mục
C3 riêng, có cùng corpus và cùng tiền xử lý. Cần nói rõ hai điều. Một, giữ lại một
từ đã bị xoá khỏi chỉ mục không phục hồi được khả năng truy xuất bằng chỉ mục đó;
fallback chỉ cứu được truy vấn mà không còn token nào. Hai, hệ thống khi đó phải
giữ cả hai chỉ mục, nên dung lượng triển khai tăng lên chứ không giảm; số liệu
mục 6.1 của Chương 4 đo riêng cả chỉ mục C4 thuần lẫn gói C4 cộng C3.

\section{Cấu trúc chỉ mục và đường đi của một truy vấn}

\begin{enumerate}
\item Chuẩn hóa văn bản và tách từ theo đối tượng tiền xử lý của M1.
\item Lọc stopword theo đúng cấu hình đã dùng khi lập chỉ mục.
\item Vector hoá truy vấn bằng bảng IDF ngoại tuyến của corpus.
\item Lấy ứng viên bằng hợp các posting list trong chỉ mục đảo.
\item Tính cosin trên tập ứng viên, chấm điểm chỉ những kết quả có điểm dương.
\item Sắp xếp giảm dần điểm; hòa điểm sắp theo \texttt{doc\_id} tăng dần để kết
quả tái lập được.
\end{enumerate}

Chỉ mục được lưu dạng thưa: posting list và vector trọng số. Hàm giải thích
(\texttt{explain}) trả về từng đóng góp vào tích vô hướng, và tổng các đóng góp
bằng đúng điểm cosin đã trả về, để người đọc có thể kiểm tra lại cách xếp hạng.

\section{Giao diện đánh giá và cổng chặn kết quả chính thức}

Bộ đánh giá dùng các định nghĩa chỉ số sau, với $K$ mặc định là 10 và có thể cấu
hình:

\begin{align}
P@K &= \frac{\text{số trúng trong Top-}K}{K}, \\
Recall@K &= \frac{\text{số trúng trong Top-}K}{\text{tổng số relevant đã confirmed}}, \\
F1@K &= \frac{2 \cdot P@K \cdot Recall@K}{P@K + Recall@K}, \\
AP &= \frac{1}{R}\sum_{i=1}^{L} P@i \cdot rel_i, \\
NDCG@K &= \frac{DCG}{IDCG},\quad DCG=\sum_{i=1}^{K}\frac{2^{rel_i}-1}{\log_2(i+1)} .
\end{align}

Ba quyết định thiết kế đáng ghi lại. Mẫu số của $P@K$ là $K$ cố định, kể cả khi
hệ thống trả về ít hơn $K$ kết quả; thiếu kết quả là sai chứ không phải là không
tính. $AP$ tính trên toàn bộ ranking có sẵn, không cắt theo $K$, nên bộ đánh giá
từ chối chạy nếu độ sâu lấy ranking nhỏ hơn số tài liệu corpus. $IDCG$ lấy từ
qrels chứ không lấy từ lần chạy, nên điểm không thể bị đẩy lên bằng cách chỉ
trả về một phần kết quả.

Query không có tài liệu relevant nào vẫn được giữ trong bình quân với giá trị 0 và
được đếm riêng, để người đọc biết chính xác MAP bị kéo xuống bao nhiêu do những
query đó. Tương tự, một truy vấn không có kết quả trả về vẫn sinh ra dòng chỉ số
bằng 0 chứ không bị loại khỏi bình quân.

\subsection{Điều kiện để coi là đánh giá chính thức}

Bộ đánh gia từ chối ghi kết quả chính thức khi có bất kỳ điều kiện nào sau đây:

\begin{itemize}
\item dưới 20 truy vấn;
\item SHA256 corpus lệch giữa manifest, qrels, tập truy vấn hoặc chỉ mục, hoặc
thiếu hẳn hash ở bất kỳ nguồn nào;
\item còn phán quyết chưa confirmed, hoặc nhãn confirmed mà điểm relevance còn
null, hoặc nhãn confirmed không ghi người gán;
\item trùng cặp (truy vấn, tài liệu), điểm relevance ngoài tập $\{0,1,2\}$, hoặc
truy vấn trong qrels không tồn tại trong tập truy vấn;
\item metadata tiền xử lý chưa nhất quán ở trạng thái official, hoặc đang dùng
tokenizer phát triển hay bộ lọc tầng 1 tạm.
\end{itemize}

SHA256 corpus cần đánh giá không bao giờ được lấy từ qrels: nếu không xác định
được từ manifest đã khóa hoặc từ chính file corpus thì kết quả chính thức bị chặn.

\subsection{Khi chưa đủ nhãn thì không xuất chỉ số}

Nếu bất kỳ truy vấn nào chưa có nhãn confirmed, bộ đánh giá ghi
\texttt{metrics\_status = unmeasured} và không ghi bảng chỉ số. Lý do rất cụ thể:
với qrels rỗng, mọi chỉ số đều bằng không một cách máy móc, và một bảng toàn số
không dễ bị đọc nhầm là hệ thống không tìm được gì.

\section{Tách kết quả kiểm thử khỏi kết quả chính thức}

Các bài kiểm thử dùng dữ liệu tạm trong thư mục làm việc của pytest. Để ngăn
tình huống số liệu kiểm thử lọt vào báo cáo, bộ đánh giá chặn ghi kết quả không
chính thức vào các thư mục kết quả chính thức; kết quả chạy thử bắt buộc ghi ra
thư mục khác, ví dụ \texttt{results/development\_eval}.
